\section{Arithmetic savings for exceptional words}\label{q:rank-witness-section}

The centered trace has been reduced to the majorants in
\Cref{q:mixed-difference}.  We now discard two further classes.
Independent lit-consistency equations provide a rank saving, while many
singleton coordinates force positive witnesses whose relations can be
summed in a triangular order.  Each saving makes its entire class
negligible before we count the remaining words.

Recurrence congruences constraining prime labels play a central role in
\cite[Sections~6--7]{HR21}.  We first prove a random-prime rank estimate
suited to the present coefficients.  The later witness argument adapts
the active-prime and triangular-constraint constructions of
\cite[Lemmas~11.6, 13.1, and~13.6]{Pilatte26}.

\subsection{Saving from independent congruences}

We first discard the words whose lit-consistency conditions contain
many independent pairs of vectors.  The point is to use rank over
$\mathbb R$, without assuming that a matrix remains invertible modulo
any of the primes being summed.  The following probability estimate
isolates the argument.

\begin{lemma}[A rank estimate for random prime divisors]
\label{q:rank-probability}
Let $\mathscr P$ be a finite nonempty set of primes, and let $\mu$ be a
probability measure on $\mathscr P$ with
$\max_{p\in\mathscr P}\mu(p)\le\alpha$.
Let $r\ge1$, let $B\in M_r(\mathbb Z)$ be nonsingular over
$\mathbb R$, and let $f:\mathscr P^r\to\mathbb Z^r$ be any function.
Suppose that $M\ge2$ and
\[
 \left|(B(y-y'))_i\right|\le M
 \quad\text{for all }y,y'\in\mathscr P^r\text{ and }1\le i\le r.
\]
If the $2r$ coordinates of $c,y\in\mathscr P^r$ are independent with
law $\mu$, then
\begin{equation}\label{q:rank-probability-bound}
 \Pbb\bigl(c_i\mid(By+f(c))_i\text{ for every }i\bigr)
 \le \left[\alpha(1+d_M)\right]^{r/2},
 \qquad d_M=\left\lfloor\frac{\log M}{\log2}\right\rfloor.
\end{equation}
\end{lemma}

\begin{proof}
Let $E(c,y)$ denote the divisibility event in the statement, and let
$y'$ be an independent copy of $y$, independent also of $c$.
Cauchy--Schwarz, applied to the conditional probability with $c$ fixed,
gives
\begin{align*}
 \Pbb(E(c,y))^2
 &=\left(\E_c\Pbb_y(E(c,y)\mid c)\right)^2\\
 &\le \E_c\Pbb_y(E(c,y)\mid c)^2\\
 &=\Pbb(E(c,y)\cap E(c,y'))\\
 &\le\Pbb\bigl(c_i\mid (B(y-y'))_i\text{ for every }i\bigr).
\end{align*}
The two copies have the same controlling primes $c_i$, so subtraction
removes the entire vector $f(c)$.

Write $v=B(y-y')$.  For a subset $Z\subseteq\{1,\ldots,r\}$ of size
$a$, the rows of $B$ indexed by $Z$ have rank $a$.  Choose an
invertible $a$-column minor in those rows.  Conditional on $y'$ and on
the coordinates of $y$ outside that minor, the equations
$v_i=0$ for $i\in Z$ determine at most one vector of values for its
$a$ remaining coordinates.  Independence and the atom bound imply
\begin{equation}\label{q:rank-zero-rows}
 \Pbb(v_i=0\text{ for all }i\in Z)\le\alpha^a.
\end{equation}
This includes $a=0$, with the empty condition having probability one.

Now fix $y,y'$.  A nonzero integer of absolute value at most $M$ has
at most $d_M$ distinct prime divisors.  Each nonzero coordinate $v_i$
therefore satisfies $c_i\mid v_i$ with conditional probability at most
$d_M\alpha$.  The controls are independent of each other and of $v$.
If the exact zero set of $v$ is $Z$, their joint conditional
probability is thus at most $(d_M\alpha)^{r-|Z|}$.
Using \eqref{q:rank-zero-rows} and summing over the possible exact zero
sets gives
\[
 \Pbb(c_i\mid v_i\text{ for every }i)
 \le\sum_{Z\subseteq\{1,\ldots,r\}}
       \alpha^{|Z|}(d_M\alpha)^{r-|Z|}
 =\left[\alpha(1+d_M)\right]^r.
\]
Taking the square root proves \eqref{q:rank-probability-bound}.
\end{proof}

We apply the lemma to one column of the prime tuples in a word of
length $2k$.  Put
\[
 r=\lfloor L^{1/50}\rfloor.
\]
Fix a column $\nu\in\{1,\ldots,J\}$ and an equality pattern for its
prime slots.  Let $\mathcal Z_\nu$ be the set of its abstract labels,
let $z_a\in\mathcal Z_\nu$ be the label at step $a$, and write $p_z$
for the numerical prime assigned to $z$.  Fix the padding values,
orientations, and the primes in every other column.  In the real
vector space with basis $(e_z)_{z\in\mathcal Z_\nu}$ define
\begin{equation}\label{q:formal-offsets}
 t_a=\frac{D_a}{p_{z_a}},\qquad
 v(i)=\sum_{a<i}t_a e_{z_a}\quad(1\le i\le2k).
\end{equation}
Each step has exactly one prime from column $\nu$.  Consequently
$t_a$ is an integer independent of every numerical prime in that
column.  Evaluating $e_z$ at $p_z$ sends $v(i)$ to the departure
offset $\xi_i$.

For two lit occurrences $i,i'$ of a label $z$, consider the pair
\begin{equation}\label{q:rank-pair}
 (e_z,v(i)-v(i')).
\end{equation}
A collection of these pairs is \emph{jointly independent} when all
its displayed vectors together are linearly independent over
$\mathbb R$.

\begin{lemma}[Discarding words with large rank]\label{q:rank}
In the sum of the nonnegative majorants from
\Cref{q:mixed-difference}, the total contribution of words having a
column with $r$ jointly independent pairs \eqref{q:rank-pair} is
\[
 O\left(\exp(-c_0 L^{203/200})\right)
\]
for some absolute $c_0>0$ and sufficiently large $L$.
The bound includes the external dimension factors in the trace
expansion and is uniform in the bin.  The threshold may depend on the
fixed parameters $W,h,l$.
\end{lemma}

\begin{proof}
Use the positive reciprocal majorant and enumeration of
\Cref{q:crude-count}.  Thus we may discard the denominators and
cutoffs, retain one reciprocal for every distinct prime label, and
include the factors $Lu(q_a)$ and $2^{S_1}$ in the crude count.  We
retain the lit-consistency conditions.  The total unrestricted count,
including external dimension factors, is
$\exp(O(L\log^2L))$.

Fix a column, its equality pattern, and the other-column and padding
data.  Choose $r$ jointly independent pairs, and denote their
controlling labels by $z_1,\ldots,z_r$ and their difference vectors
by $w_1,\ldots,w_r$.  The controlling labels are distinct.  Modulo
$C=\operatorname{span}(e_{z_1},\ldots,e_{z_r})$, the images of the
$w_i$ remain independent: a dependence would express a nontrivial
linear combination of the $w_i$ as a combination of the $e_{z_i}$,
contrary to joint independence.  Therefore the row matrix of the
$w_i$, after deleting all controlling columns, has rank $r$.
Choose an invertible $r$-column minor $B_0$ of that matrix.

The entries of the entire row matrix, and hence the choice of a
minor by any fixed deterministic rule, are independent of all the
numerical primes in column $\nu$.  Fix the prime variables in that
column outside the controlling coordinates and the minor.  Write
$c_i=p_{z_i}$ for the controls and $y$ for the vector of primes in
the minor.  The selected lit-consistency conditions become
\begin{equation}\label{q:rank-congruences}
 c_i\mid (B_0y+A_0c+d_0)_i\qquad(1\le i\le r),
\end{equation}
with fixed integer matrices $B_0,A_0$ and a fixed integer vector
$d_0$.  No invertibility assertion modulo $c_i$ is being made.

Enlarge the positive sum by dropping numerical distinctness among
abstract labels, as well as bin and closure restrictions.  Continue
to attach one factor $1/p$ to each abstract label: if two coordinates
now take the same prime, their product weight is $1/p^2$.
The enlarged reciprocal sum is consequently a product sum.  Its
normalized coordinates have independent law
\[
 \mu_\nu(p)=\frac{1}{V_\nu p}\quad(p\in P_\nu),
 \qquad \max_p\mu_\nu(p)\le\frac{1}{V_\nu H_0}
 \le\frac1{H_0}.
\]
This independence allows numerical coincidences.  In particular, after
subtraction in \Cref{q:rank-probability}, the controls remain
independent of $B_0(y-y')$, even on a sample space that permits such
coincidences.

We check the size hypothesis after these enlargements.  The fixed
padding values came from eligible steps, so $q_a\le e^{100L+1}$.
The product of the primes in the other columns is at most $e^{2L}$
by the geometric spacing of their bands.  Hence
\[
 |t_a|\le h e^{102L+1}.
\]
Each row of the difference matrix has sum of absolute coefficients
at most $2kh e^{102L+1}$.  Since every prime in the resampled column
is at most $e^L$, we obtain
\[
 |(B_0(y-y'))_i|\le 2kh e^{103L+1}=e^{O(L)}.
\]
For fixed $h$ and sufficiently large $L$, the constant in this last
exponent is absolute.  In \Cref{q:rank-probability} we may therefore
take $d_M=O(L)$, $\alpha=(V_\nu H_0)^{-1}$, and
$f(c)=A_0c+d_0$.  The probability of \eqref{q:rank-congruences} is at
most
\[
 \left(\frac{O(L)}{V_\nu H_0}\right)^{r/2}
 \le \exp(O(r\log L))H_0^{-r/2}.
\]
Multiplying by the unrestricted harmonic masses restores the
unnormalized reciprocal sum with this same relative saving.
The estimate is uniform in every fixed outside variable, so these
variables may now be summed.

The column, selected pairs and selected minor have at most
$J(2k)^{O(r)}$ descriptions.  Their cost is absorbed in
$\exp(O(r\log L))$.  Together with the crude enumeration, the total
is bounded by
\[
 \exp\left(O(L\log^2L)+O(r\log L)-\frac r2\log H_0\right).
\]
Since $\log H_0=L^{199/200}$ and
$r=L^{1/50}+O(1)$, its negative term is
$(\tfrac12+o(1))L^{203/200}$ and dominates both positive terms.
This proves the lemma.
\end{proof}

We may therefore restrict the remaining trace sum to words for which,
in every column, a maximal jointly independent collection of pairs
\eqref{q:rank-pair} has fewer than $r$ members.  This quantitative
restriction will constrain the equality patterns once the words with
many singleton labels have also been discarded.

\subsection{Singleton primes and positive witnesses}

We next dispose of words having many singleton primes.  Without the
vertex deletion, integration of any singleton centered factor would give
zero.  The mixed difference in \Cref{q:mixed-difference} measures the
failure of this cancellation caused by deletion.  We show that a nonzero
mixed difference forces many positive witnesses, and that these witnesses
supply successively usable congruences on distinct prime labels.

Put
\[
 T=\lfloor L^{1/12}\rfloor.
\]
\begin{lemma}[A common configuration of witnesses]\label{q:witness-cover}
Fix a numerical main word with singleton set $\mathcal U$, where
$S_1=|\mathcal U|>L^{1/4}$, and fix a draw of the residue coordinates.
If $\Delta G\ne0$, then some hybrid contains $T$ positive witnesses
attached at main departures, each of length at most $s$, with distinct
marked primes $y_1,\ldots,y_T\in\mathcal U$.  The prime $y_j$ occurs
in witness $j$ and in none of the other $T-1$ witnesses.  The combined
length of the main word and these witnesses is at most $4L$.
\end{lemma}

\begin{proof}
For the fixed numerical word, let $\mathscr W$ be the finite set of all
possible witness tests at its departure sites.  Write $I_w$ for the
indicator that the preselected numerical word $w$ is positive at its
assigned departure.  Then
\[
 G=\prod_{w\in\mathscr W}(1-I_w).
\]
A repeated test in this product is harmless because its factors take
values in $\{0,1\}$.  Expand this finite product algebraically and apply
the mixed difference in $\mathcal U$.  If the union of the prime supports
of a subfamily $\mathscr A\subseteq\mathscr W$ omits a prime
$p\in\mathcal U$, then $\prod_{w\in\mathscr A}I_w$ is independent of
the $p$ coordinate.  Its full mixed difference is therefore zero.
Consequently $\Delta G\ne0$ implies that some subfamily whose supports
cover $\mathcal U$ has a nonzero mixed difference.  Its intersection
indicator is then one in at least one hybrid.  All witnesses in that
subfamily are positive in this single configuration.

Choose an inclusion-minimal subfamily covering $\mathcal U$ there.
Each member contains a prime of $\mathcal U$ occurring in no other
member, since otherwise that member could be removed.  A witness has at
most $sJ$ distinct $P$ labels, so this subfamily has at least
$S_1/(sJ)$ members.  For sufficiently large $L$,
\[
 \frac{S_1}{sJ}>\frac{L^{3/20}}{J}>T.
\]
Keep any $T$ members and one private prime in each.  Their private
primes remain private in the selected subfamily.  Finally,
\[
 2k+sT\le 2L+L^{11/60}\le4L.
\]
The algebraic expansion above has been used only to deduce existence.
No bound is obtained by summing the absolute values of its terms, and
no enumeration of all its subfamilies is charged.
\end{proof}

The cover gives each selected witness a private singleton prime, but
privacy alone does not give a congruence with an invertible coefficient
in that prime.  We first identify labels having such a coefficient in
an internal relation.  For a minimal prohibited word with steps
$E_1,\ldots,E_m$, call a $P$ label $x$ \emph{active}
if there are a $P$ label $c$ of the same word and an interval
$I\subseteq\{1,\ldots,m\}$ such that
\begin{equation}\label{q:active-definition}
 c\mid\sum_{a\in I}E_a,
 \qquad
 \sum_{\substack{a\in I\\x\mid E_a}}E_a\not\equiv0\pmod c.
\end{equation}
In particular $x\ne c$.  With every other numerical label fixed, the
first sum is linear in $x$ and its coefficient is invertible modulo
$c$.  These are exactly the conditions needed for a reciprocal-prime
saving.

\begin{lemma}[Active labels in a minimal witness]\label{q:active}
Let $E_1,\ldots,E_m$ be a minimal prohibited word.  For every $P$ label
$y$ in this word, either $y$ is active or there are an active label $z$
and a step $i$ containing $z$ such that
\begin{equation}\label{q:active-prefix}
 \sum_{\substack{a<i\\y\mid E_a}}E_a\not\equiv0\pmod z.
\end{equation}
The relations establishing activity have the form
\eqref{q:active-definition} with contiguous intervals.
\end{lemma}

\begin{proof}
Consecutive tuples differ, and each label occupies an interval of step
indices.  Thus some label $z$ enters for the first time on the last
step.  Let $c$ be the first-step label controlling the prohibited suffix
of the word.  Since $c$ is absent from the last tuple and divides neither
$h$ nor a padding factor, $c\nmid E_m$.  The contribution of $z$ to the
prohibited suffix is exactly $E_m$, so $z$ is active.  The same argument
shows that every label used only on the last step is active.

Now suppose that $y$ occurs before the last step.  Restricting attention
to steps $1,\ldots,m-1$, its occurrences form a nonempty interval
$[a,b]$, where $b<m$; $y$ may also occur on step $m$.
If $E_a+\cdots+E_b$ is nonzero modulo $z$, then
\eqref{q:active-prefix} holds at $i=m$.  Otherwise
\begin{equation}\label{q:active-zero-prefix}
 z\mid E_a+\cdots+E_b.
\end{equation}
Every step before the last is nonzero modulo $z$, so $b>a$.
Reverse the contiguous subword $E_a,\ldots,E_m$, negating all its steps.
Its first step contains $z$, its last does not, and its suffix consisting
of the reversed steps $E_b,\ldots,E_a$ has sum divisible by $z$.
This suffix starts strictly after the first step because $b<m$, and
strictly before the last step because $b>a$.  The interval appearances
and unequal consecutive tuples persist under reversal.  The reversed
subword is therefore forward prohibited.  Minimality forces $a=1$.

Choose a label $z'$ newly entering on step $2$, which is possible because
the first two tuples differ.  Let $b'$ be the smaller of $b$ and the last
index carrying $z'$.  Its contribution to \eqref{q:active-zero-prefix}
is $E_2+\cdots+E_{b'}$.  If that contribution were zero modulo $z$,
then $b'>2$, since $z\nmid E_2$.  The reversed subword
$E_m,\ldots,E_2$ would then be forward prohibited, with controlling
prime $z$ and suffix $E_{b'},\ldots,E_2$.  It is shorter than the
original word, contradicting minimality.  Hence $z'$ is active through
\eqref{q:active-zero-prefix}.  Since $y$ occurs on step $1$ and $z'$ does
not, its contribution before step $2$ is $E_1\not\equiv0\pmod{z'}$.
This proves \eqref{q:active-prefix} with $i=2$ and $z'$ in place of $z$.
\end{proof}

We now have two ways to obtain a usable congruence.  An active label
already appears with an invertible coefficient in an internal relation.
For a private label that is not active, Lemma~\ref{q:active} supplies a
nonzero prefix contribution modulo another active label.  If that active
label also occurs in another witness, positivity lets us compare the
two departures to obtain a congruence.  The next proof either selects
enough internal relations or arranges enough of these comparisons in
an order suitable for elimination.

\begin{lemma}[Elimination of the singleton class]\label{q:singletons}
For a fixed bin, the total contribution to the absolute majorants in
\Cref{q:mixed-difference} from main words with $S_1>L^{1/4}$, including
the external factor $|\mathcal D|^2M$ in the trace expansion, is at most
\[
 \exp(-L^{21/20})
\]
for sufficiently large $L$, with the fixed parameters held fixed.
\end{lemma}

\begin{proof}
Whenever the integrand is nonzero, Lemma~\ref{q:witness-cover} supplies
$T$ witnesses positive in one common hybrid, with marked private
singletons.  Order them as $w_1,\ldots,w_T$ by their attachment indices
$a_1\le\cdots\le a_T$ on the main word; order ties arbitrarily.
Write $t(y)$ for the unique main step carrying a main singleton $y$.
Thus the main position of $y$ has been passed at attachment $a$ exactly
when $t(y)<a$.
We will select at least $\lfloor T/8\rfloor$ prime variables and order
their relations so that each relation is independent of all later
selected variables.  In the first two cases we select active labels,
which need not be the marked singletons.  In the remaining case we
select the private marked singletons themselves.

\emph{Internal relations.}
Suppose at least $T/8$ witnesses have an active label absent from every
earlier witness.  Choose one such label from each of these witnesses,
and record an internal relation establishing its activity.  The chosen
labels are distinct.  Each earlier relation involves only labels in its
own witness, so it is independent of every later chosen label, including
as a controlling modulus.  Sum the chosen prime values in reverse
witness order.  When one is summed, all earlier relations are independent
of that variable, while its own relation has an invertible coefficient.
This gives at least $\lfloor T/8\rfloor$ successive reciprocal-prime
savings.  If instead at least $T/8$ witnesses have an active label absent
from every later witness, apply the same argument with the order
reversed.

\emph{Comparison relations.}
Assume neither alternative holds.  Apart from fewer than $T/4$
witnesses, every active label occurs in both an earlier and a later
witness.  In each remaining witness $w_j$, its private marked singleton
$y_j$ is not active.  By Lemma~\ref{q:active}, choose an active label
$z_j$ and a departure within $w_j$ whose prefix has nonzero
$y_j$-contribution modulo $z_j$.

First consider those $j$ for which $t(y_j)\ge a_j$.  Choose an earlier
witness $w_h$, $h<j$, containing $z_j$, and a departure there using
$z_j$.  Let $\beta_j$ and $\beta_h$ be the offsets of these selected
departures from the respective starts of $w_j$ and $w_h$.
Positivity in the common hybrid gives
\begin{equation}\label{q:earlier-comparison}
 z_j\mid(\xi_{a_j}+\beta_j)-(\xi_{a_h}+\beta_h)
      =\sum_{a_h\le v<a_j}D_v+\beta_j-\beta_h.
\end{equation}
The main segment omits $y_j$ because $t(y_j)\ge a_j$, and $w_h$ omits
it by privacy.  Its only contribution is therefore through $\beta_j$,
where it has the nonzero prefix
contribution given by Lemma~\ref{q:active}.

For two selected witnesses $w_j,w_i$ with $j<i$, the private label
$y_i$ occurs in neither witness prefix in
\eqref{q:earlier-comparison}.  It is also absent from its main segment,
since
\[
 t(y_i)\ge a_i\ge a_j.
\]
Thus the comparison belonging to $w_j$ is independent of every later
selected variable $y_i$.  This includes its controlling prime: $z_j$
occurs in both $w_j$ and $w_h$, whereas each selected $y_i$ is private.
Eliminating the selected labels in decreasing attachment order is
therefore valid.  Equal attachment indices cause no difficulty, as the
main segment still ends strictly before that index.

For the witnesses with $t(y_j)<a_j$, choose instead a later witness
containing $z_j$.  The comparison uses the main segment starting at
$a_j$ and ending just before that later attachment, so it omits $y_j$.
Order this group by decreasing attachment index.  A selected label from
a smaller attachment has its unique main position still smaller than
that attachment, and hence lies before the start of every earlier
comparison segment in this order.  Privacy removes it from the two
witness prefixes as well.  The same reverse-elimination argument applies.
One of these two groups has at least $T/8$ witnesses for sufficiently
large $T$: there are more than $3T/4$ candidates before dividing them
according to whether their main position has passed.

\emph{Summing the relations.}
In all cases, record the chosen relations, their order, their controlling
labels, and the nonvanishing coefficient tests.  Fix the equality pattern,
the numerical padding factors, and all nonselected numerical labels.
Every selected congruence is linear in its selected prime $x$; its
coefficient is a unit modulo its controlling prime $p$.  The latter is
unselected or belongs to an earlier relation in the chosen order.
The reciprocal sum is bounded by
\[
 \sum_{\substack{H_0\le x\le e^L\\x\equiv c\pmod p}}\frac1x
 \le\frac1{H_0}+\frac Lp
 \le\frac{1+L}{H_0},
\]
where enlarging from primes to integers only increases the sum.  Retain
the nonvanishing coefficient condition during each elimination; if it
fails for the remaining fixed data, that inner sum is zero.  After
summing a selected variable, drop its already used relation.  The explicit
dependencies above ensure that all still retained relations are
independent of the variable just eliminated.

It remains to justify the residue cost and the number of records.
Bound $|\Delta G|$ by $2^{S_1}$, drop denominators and cutoffs from $R$,
and retain its padding divisibilities.  The factor $b_{\mathcal L}$
already supplies one reciprocal for every main $P$ label.  Each new
witness $P$ label and each distinct main or witness $Q$ label has a
positivity test in its original residue coordinate, which no hybrid
changes.  Their joint average supplies one reciprocal per distinct
such label.  Constraints on overwritten main coordinates may be dropped.
In particular no choice of hybrid needs to be counted: these retained
residue tests are unchanged throughout the hybrids, and the recorded
numerical relations are necessary consequences of positivity in the
common one.

The total recorded length is at most $4L$.  Lemma~\ref{q:crude-count}
therefore bounds the equality patterns, numerical reciprocal sums before
the selected savings, signs, and weights by $\exp(O(L\log^2L))$.
Attachments, marked labels, internal interval endpoints, partner-witness
indices, and elimination orders have only polynomially many options per
recorded slot, so their inclusion preserves this bound.  The factors
$2^{S_1}$ and $|\mathcal D|^2M=\exp(O(L))$ do so as well.  We have
proved the bound
\[
 \exp(O(L\log^2L))
 \left(\frac{1+L}{H_0}\right)^{\lfloor T/8\rfloor}.
\]
Finally, $T\log H_0\asymp L^{1/12+199/200}=L^{647/600}$.
This dominates both $L\log^2L$ and $L^{21/20}$, proving the assertion.
\end{proof}
